\documentclass[a4paper,12pt, french]{article}

\usepackage{layout_evaluation}

\classe{Terminale STMG}
\titre{Suites arithmétiques}
\evalnumber{}
\evaltype{DS}
\date{}

\begin{document}

\name

\emph{Durée : 1H\qquad Calculatrice interdite \qquad \textbf{Les exercices sont indépendants.}}

\vspace{-3mm}

\subsection*{Exercice 1 (8 points):}
On s'intéresse au nombre de voitures produites dans une usine Peugeot. En 2020 l'usine a produit 15\ 000 voitures. L'entreprise entend augmenter sa production de 200 voitures par an. On note $u_n$ le nombre de voitures produites à l'année $2020+n$.

\begin{enumerate}
    \item A quoi correspond $u_0$ ? Donner sa valeur.
    \item Donner la relation entre $u_{n+1}$ et $u_n$.
    \item Donner l'expression du terme général $u_n$.
    \item En déduire le nombre de voitures produites en 2026 ?
    \item En quelle année la production de voiture atteindra 20\ 000 en une année ?
\end{enumerate}

\subsection*{Exercice 2 (6 points):}
Soit la suite $(v_n)$ définie par $v_{n+1} = v_n-12$ et $v_1=4$.

\begin{enumerate}
    \item Donner la raison de la suite et déterminer les quatre premiers termes.
    \item Quelle est la moyenne arithmétique de $v_1$ et $v_4$ ?
    \item A l'aide du symbole de somme $\Sigma$, écrire la somme de ces quatre premiers termes.
    \item Calculer cette somme.
\end{enumerate}

\subsection*{Exercice 3 (2 points):}
Calculer $\sum_{k=0}^{3} (k^2+1)$

\subsection*{Exercice 4 (4 points+1 point bonus):}
Soit une suite arithmétique de premier terme $w_1 = 100$ et de raison $r=2$.
\begin{enumerate}
    \item Que vaut $w_{9}$ ?
    \item Donner la formule permettant de calculer la somme des 9 premiers termes.
    \item (Bonus) Effectuer le calcul.
\end{enumerate}

% \subsection*{Exercice bonus \diff (1 point):}
%     Soit la suite $(u_n)$ définie par $u_0=-154$ et sa raison $r=13$. Calculer le terme $u_{11}$.

\end{document}
